\begin{textsample}{POS Dim 5 – global\_north\_2025\_06 – Score 14.00 – global\_north\_2025\_06/125816992.txt}  \label{ex:f5_pos_085}
From Durban to Geneva : How the global human rights industry turned on Israel The hostile takeover of key humanitarian organisations has been visible for more than two decades June 12, 2025 14:35 UNhinged : the 2001 Durban conference kicked off a global campaign to isolate and demonise the Jewish state ( Image : Getty ) 3 min read When \textbf{Amnesty} \textbf{International} and Human Rights Watch were founded in 1961 and 1978 respectively -- both by Jews and Zionists -- they quickly earned reputations as principled defenders of universal human rights.
Yet over time, both organisations have drifted from their original mission of confronting the world ' s most brutal regimes.
Today, they are increasingly politicised, with a marked and obsessive hostility towards Israel.
The hostile takeover became clearly visible in August 2001, when the NGO Forum of the UN ' s Conference on Racism brought 5,000 activists from self-proclaimed human rights groups to Durban, South Africa.
The orchestrated assemblage declared Israel to be guilty of apartheid, \textbf{genocide}, colonialism, among similar propaganda @ @ @ @ @ @ @ @ @ @, boycott campaigns and other forms of demonisation based on exploiting the principles and frameworks of human rights.
Twenty-two years later, immediately after the October 7 atrocities, the world-wide propaganda attacks ("the 8th front"of the war ) highlighted the same slogans in much more virulent form, feeding blood libels, antisemitic violence and intimidation.
The failure of the Israeli government, including the IDF, as well as the leaders of major Jewish organisations, to \textbf{recognise} and prioritise systematic responses to NGO warfare allowed this danger to fester and expand.
The malign political influence of groups such as \textbf{Amnesty} \textbf{International}, Human Rights Watch and their numerous allies increased continuously.
But the IDF and various ministries paid little attention to their propaganda reports, parroted in headline articles by prominent journalists around the world, which labelled every response to mass terror as a"war \textbf{crime}".
In 2009, the Goldstone report ( the UN"Fact Finding Mission on the Gaza Conflict") accused Israel of"possible \textbf{crimes} against \textbf{humanity} @ @ @ @ @ @ @ @ @ @ created \textbf{International} Criminal \textbf{Court}. \textbf{Amnesty} \textbf{International} wrote the list of alleged war \textbf{crimes}, and the majority of the more than 500 citations of"evidence"in the final document were sourced to 50 anti-Israel NGOs.
Months later, after \textbf{Judge} Richard Goldstone met with critics ( including myself ), he acknowledged that his document was deeply biased and inaccurate, but the damage was done.
The threat of \textbf{international} \textbf{legal} action against soldiers got the attention of the IDF, government lawyers and other officials, but the responses were ad hoc.
The counter-strategy consisted of claims that the IDF was"the world ' s most moral army", that Israel investigated all \textbf{allegations} of \textbf{violations}, as well as numerous learned \textbf{legal} briefs arguing that the ICC and other \textbf{international} frameworks lacked \textbf{jurisdiction}.
This approach had little to no impact on the lawfare and propaganda campaigns that singled out Israel for demonisation.
On the contrary, the advocacy NGOs and their allies in the media, UN, and university campuses ( particularly under the headings of @ @ @ @ @ @ @ @ @ @ disproportionate attacks, and their influence increased with every round of the Gaza conflict.
On social media platforms and Wikipedia, NGOs moved quickly to dominate the content on Israel, Palestinian victimhood and human rights -- facilitated by public relations consultants and enabled through massive budgets ( \textbf{Amnesty} reported a budget of? 380? million in 2023 ; HRW ' s is \$100? million ).
In addition to the NGO superpowers, at least 300 smaller groups operate in different venues to supplement and amplify the campaigns against Israel, including through major church organisations and sub-networks of Palestinian and Israeli advocacy organisations largely funded by"civil society"grants from European government supporters.
Organisations such as Al Haq and the Gaza-based Palestinian Centre for Human Rights ( PCHR ) obtained UN recognition, enabling them to present \textbf{allegations} and reports against Israel at every session of the Human Rights Council in Geneva, where diplomats applaud and journalists quote their \textbf{accusations}.
In October, 2021, the Israeli Ministry of Defence belatedly designated Al-Haq and five related NGOs as @ @ @ @ @ @ @ @ @ @ of organisations"operated on behalf of the Popular Front for the Liberation of Palestine ( PFLP ), which is banned in the US, Canada, the EU and elsewhere.
But government follow-up was sporadic, and, as of June 2025, officials from Al Haq and the others continue to lobby European ministries and UN delegations, and file lawsuits seeking to block defence \textbf{exports} to Israel through"war \textbf{crimes}"\textbf{allegations}.
Throughout this period, different Israeli government ministries fought turf battles over budgets and areas of responsibility.
At the same time, right-wing politicians focused narrowly on the Israeli political NGOs affiliated with the New Israel Fund and introduced legislation seeking to block their funding from European governments.
These efforts did not lead to regulatory changes, and instead increased their external influence in normalising the abuse of human rights principles for demonising Israel, including the campaign based on the Durban strategy of uniquely labelling Israel as inherently guilty of apartheid.
In reviewing the history of NGO warfare, the processes leading to the current @ @ @ @ @ @ @ @ @ @ \textbf{genocide} and deliberate starvation, are all too visible.
And although many years late, effective counter strategies in this important 8th front of Israel ' s war must focus on the sources of NGO power -- specifically their donors and enablers.
In recent months, former senior officials of HRW, \textbf{Amnesty}, and Doctors Without Borders have stepped forward and condemned the widespread antisemitism and anti-Israel obsession within their organisations.
This is an important step, although far too late to undo the damage.?

% matched lemmas: accusation, allegation, amnesty, court, crime, export, genocide, humanity, international, judge, jurisdiction, legal, recognise, violation
\end{textsample}
